\subsection{The actual cubic-plethysm multiplicity algebra}

\begin{proposition}[PLE-F8: the obstructing ray is the binary-dihedral algebra]
\label{prop:actualplethysmalgebra}
Let $V=\mathbb C^2$, let $W=\operatorname{Sym}^3V$, and write
$S=\operatorname{Sym}(W)=\mathbb C[a,b,c,d]$ using the basis
$a=x^3$, $b=x^2y$, $c=xy^2$, $d=y^3$ of $W$.
Here $a,b,c,d$ are algebraically independent generators of
$\operatorname{Sym}(W)$; the monomials in $x,y$ label the basis
vectors of $W$.
The torus $\operatorname{diag}(t,t^{-1})$ assigns these variables
weights $3,1,-1,-3$, and the raising operator is the derivation
\[
 E=a\frac{\partial}{\partial b}
      +2b\frac{\partial}{\partial c}
      +3c\frac{\partial}{\partial d}.
\]
Define the graded algebra
\[
 T=\bigoplus_{s\ge0}T_s,\qquad
 T_s=\ker(E)\cap(S_{4s})_{2s},
 \tag{PLE-F8a}
\]
where the first subscript on $S$ is ordinary polynomial degree and
the second is torus weight. Then $T_s$ is the highest-weight
multiplicity space for
$S^{(7s,5s)}V$ in $\operatorname{Sym}^{4s}(\operatorname{Sym}^3V)$,
and there is an explicit graded-algebra isomorphism
\[
 T\simeq
 \mathbb C[U,V_0,Z]/(Z^2-UV_0^2-4U^4),
 \qquad \deg(U,V_0,Z)=(2,3,4).
 \tag{PLE-F8b}
\]
Thus the ring in PLE-F5 is the actual characteristic-zero plethysm
multiplicity algebra on this ray, rather than merely another algebra
with the same Hilbert function.

The algebra has a homogeneous basis
$\{U^iV_0^jZ^\epsilon:i,j\ge0,\ \epsilon\in\{0,1\}\}$
with nonnegative integral multiplication constants. It admits the
flat integral model
\[
 \mathcal T=\mathbb Z[U,V_0,Z]/(Z^2-UV_0^2-4U^4).
 \tag{PLE-F8c}
\]
For every finite field $\mathbb F_q$ with $q\equiv1\pmod{12}$
and prime $\ell\nmid6q$, PLE-F7 computes the ordinary and
intersection cohomology of its fiber: the only nonzero ordinary
compactly supported group is $H_c^4=\mathbb Q_\ell(-2)$, and
$IC=\mathbb Q_\ell[2]$. Nevertheless its Hilbert function is the
Kahle--Micha\l ek ray, which is not a homogeneous rational Ehrhart
function.
\end{proposition}

\begin{proof}
The action $E(x)=0$, $E(y)=x$ on $V$ gives exactly the displayed
derivation on $S$. Multiplication adds ordinary degrees and torus
weights, and the kernel of a derivation is a subalgebra, so (PLE-F8a)
is an algebra. Each irreducible polynomial $GL_2$ representation of
highest weight $(p,q)$ is $\det^q\otimes\operatorname{Sym}^{p-q}V$.
In its monomial weight basis the raising operator kills exactly its
highest vector, of torus weight $p-q$. All constituents of $S_{4s}$
have central weight $12s$, since a central scalar acts on $W$ with
weight three. Consequently a highest vector of torus weight $2s$
has $p+q=12s$, $p-q=2s$, hence $(p,q)=(7s,5s)$.
Complete reducibility in characteristic zero now identifies the
dimension and the space in (PLE-F8a) with the stated multiplicity
space, after choosing the usual highest line in each irreducible.
This identification respects the canonical Cartan multiplication.
Choose compatible highest vectors
$v_s=x^{2s}\otimes(x\wedge y)^{\otimes5s}$ in
$S^{(7s,5s)}V\simeq\operatorname{Sym}^{2s}V\otimes(\det V)^{5s}$.
Evaluation $F\mapsto F(v_s)$ identifies
$\operatorname{Hom}_{GL_2}(S^{(7s,5s)}V,S_{4s})$ with $T_s$.
Normalize the Cartan injection
$S^{(7(s+t),5(s+t))}V\to
 S^{(7s,5s)}V\otimes S^{(7t,5t)}V$
to send $v_{s+t}$ to $v_s\otimes v_t$.
Composing it with $F\otimes G$ and multiplication in $S$ sends
$v_{s+t}$ to $F(v_s)G(v_t)$. Hence evaluation is an isomorphism
of graded algebras for the canonical multiplication of GCT VI,
Section 4.1.2, equation (4.5), printed p.~69, constructed in
equations (4.2)--(4.3), printed p.~67.
For the $\mathfrak{sl}_2$ representation facts, see
Borcherds--Haiman--Reshetikhin--Serganova--Johnson-Freyd,
\emph{Lectures in Lie Groups and Quantum Groups}, Part I,
Example 5.2.0.5 and Proposition 5.2.0.7, printed pp.~68--69,
and Theorem 7.6.1.3, printed p.~167; the passage to $GL_2$ here
adds the scalar central character just computed.
Kahle--Micha\l ek, Theorem 5, therefore gives for $s\ge1$
\[
 \dim T_s=\phi(s)=\left\lfloor\frac{s}{3}\right\rfloor
                      +\mathbf1_{\{2\mid s\}}.
 \tag{PLE-F8d}
\]
Their notation $(7s,5s,0)$ is the same partition with a trailing
zero. Specializing the Schur-character identity to two variables
preserves this coefficient: Schur polynomials with more than two
rows vanish, and those with at most two rows are linearly independent.
At $s=0$, both sides are one directly.

We construct the isomorphism. Put
\[
 h=ac-b^2,\qquad
 j=a^2d-3abc+2b^3,
\]
and
\[
 \delta=a^2d^2-6abcd+4ac^3+4b^3d-3b^2c^2.
\]
Direct application of $E$ gives $E(a)=E(h)=E(j)=0$, and polynomial
expansion gives the identity
\[
 j^2+4h^3=a^2\delta.
 \tag{PLE-F8e}
\]
Applying $E$ to this identity proves $a^2E(\delta)=0$; since
$S$ is a domain, $E(\delta)=0$.
The pairs (ordinary degree, torus weight) of $a,h,j,\delta$ are
respectively $(1,3),(2,2),(3,3),(4,0)$.
Define actual elements of $S$ by
\[
 u=h^2\delta,\qquad
 v=aj\delta^2,\qquad
 w=a^2h\delta^3,\qquad z=w-2u^2.
\]
The pairs for $u,v,z$ are $(8,4),(12,6),(16,8)$, so they belong
to $T_2,T_3,T_4$, respectively. From (PLE-F8e) one obtains
\[
 v^2+4uw=a^4\delta^5,\qquad
 w^2=u(v^2+4uw).
\]
It follows that
\[
 z^2=(w-2u^2)^2=uv^2+4u^4.
 \tag{PLE-F8f}
\]
Thus $U\mapsto u$, $V_0\mapsto v$, $Z\mapsto z$ defines a graded
homomorphism from the right side of (PLE-F8b) into $T$.

We prove its injectivity, including the algebraic independence needed
for that step. Suppose $F(u,v)=0$ for a nonzero polynomial
$F(X,Y)\in\mathbb C[X,Y]$. Remove the largest common power of $X$
from $F$. Since $u\ne0$ and $S$ is a domain, the remaining polynomial
still gives a relation, and now $F(0,Y)\ne0$.
For instance, $u=4$ at $a=c=1$, $b=d=0$, verifying $u\ne0$.
Specialize $a=1$, $b=c=0$. Then $h=0$, $j=d$, $\delta=d^2$,
so $u=0$ and $v=d^5$. The relation would give
$F(0,d^5)=0$ in $\mathbb C[d]$, impossible for a nonzero
$F(0,Y)$. Hence $u,v$ are algebraically independent.

In the field $K=\mathbb C(U,V_0)$, the element
$UV_0^2+4U^4=U(V_0^2+4U^3)$ has valuation one at the prime $U$.
It is therefore not a square, since every square has even valuation.
Consequently $Z^2-UV_0^2-4U^4$ is irreducible over $K$.
Algebraic independence embeds $K$ in $\operatorname{Frac}(S)$ via
$U,V_0\mapsto u,v$, and the element $z$ satisfies this irreducible
quadratic. The resulting map from the degree-two field extension of
$K$ is injective. Restricting to its polynomial subalgebra proves
injectivity of the homomorphism in (PLE-F8b). Equivalently, no
relation $A(u,v)+zB(u,v)=0$ is possible: if $B\ne0$, it would put
$z$ in $K$ and make the displayed element a square, and if $B=0$,
algebraic independence forces $A=0$.

The source of this homomorphism is free over $\mathbb C[U,V_0]$
with basis $1,Z$, because its relation is monic quadratic in $Z$.
Its Hilbert series is
\[
 \frac{1+t^4}{(1-t^2)(1-t^3)}.
\]
PLE-F5 computed its coefficient in degree $s$ as $\phi(s)$.
By (PLE-F8d), the injective homomorphism is therefore an isomorphism
on each finite-dimensional graded piece, hence on the graded algebra.
This proves (PLE-F8b) without presupposing a presentation theorem
for the full algebra of binary-cubic covariants.

The same monic relation shows that the monomials stated in the
proposition form a basis over $\mathbb Z$ for (PLE-F8c).
They span freely over every coefficient ring. In multiplying two
such basis monomials, the only needed reduction is
$Z^2=UV_0^2+4U^4$, whose coefficients are nonnegative integers.
Thus the multiplication constants have the claimed positivity.
Freeness over $\mathbb Z$ also proves flatness of $\mathcal T$,
and all its fibers have the same graded Hilbert function.
Its complex fiber is $T$ by the proved isomorphism. The finite-field
statement is exactly PLE-F7 for the identical equation and grading;
the non-Ehrhart assertion is Kahle--Micha\l ek, Theorem 5 and
Remark 6, as used in P-F1 and PLE-F6. This completes the proof.
\end{proof}

\paragraph{What passes to finite characteristic, and what does not.}
The finite-field scheme here is an explicit flat model of the
characteristic-zero multiplicity algebra. It is not asserted to be
the kernel of the same Lie-algebra raising operator after reduction
modulo every prime. In characteristic $p$, a derivation kills
$p$th powers, so that kernel can acquire additional elements and
cannot be substituted for highest-weight multiplicities without
further arguments. The characteristic-zero isomorphism and the
finite-field geometric model are the claims proved above.

This result removes the possible objection that the binary-dihedral
example only reproduced the numerical Hilbert function of an
unrelated algebra. On this actual plethysm multiplicity algebra,
normality, a positive multiplicative basis, and a pure Tate model
still do not force the literal homogeneous PH1 conclusion.
The coordinate-ring degree and the cohomological degree remain
different gradings. The richer compatibility and polyhedral
requirements of a canonical-basis or crystal construction have not
been proved or refuted by this basis alone. Nor does this example
refute the narrow quantitative PH2 or PH3 bounds: it already has
the displayed positive rational form with denominator weights two
and three. No historical priority is asserted for these explicit
covariant formulas or their algebraic consequences.
